\begin{table}[H]
\centering
\footnotesize
\setlength{\tabcolsep}{3.5pt}
\begin{tabular}{lccc@{\hspace{9pt}}ccc}
\toprule
 & \multicolumn{3}{c}{CIFAR-100 ($K{=}20$ coarse)} & \multicolumn{3}{c}{ImageNet ($K{=}30$ WordNet)} \\
\cmidrule(lr){2-4}\cmidrule(lr){5-7}
model & real & star & depth ($z$) & real & star & depth ($z$) \\
\midrule
ViT-T & $-0.024$ & $-0.030$ & $+0.006$ (+0.4) & $-0.007$ & $-0.006$ & $-0.001$ (-0.1) \\
ViT-S & $-0.007$ & $-0.044$ & $+0.036$ (+2.4) & $-0.017$ & $-0.008$ & $-0.009$ (-1.0) \\
ViT-B & $+0.004$ & $-0.056$ & $+0.060$ (+4.5) & $-0.024$ & $-0.013$ & $-0.011$ (-0.9) \\
ViT-L & $-0.016$ & $-0.056$ & $+0.040$ (+4.1) & $-0.025$ & $-0.013$ & $-0.012$ (-1.3) \\
DINO-B & $-0.036$ & $-0.043$ & $+0.007$ (+0.9) & $-0.017$ & $-0.006$ & $-0.010$ (-1.6) \\
DINOv2-S & $-0.020$ & $-0.027$ & $+0.008$ (+0.7) & $-0.004$ & $-0.001$ & $-0.003$ (-0.2) \\
DINOv2-B & $+0.012$ & $-0.027$ & $+0.039$ (+5.2) & $-0.010$ & $-0.002$ & $-0.008$ (-0.8) \\
DINOv2-L & $+0.015$ & $-0.027$ & $+0.042$ (+7.0) & $-0.033$ & $-0.004$ & $-0.029$ (-2.5) \\
DINOv2-G & $-0.003$ & $-0.030$ & $+0.027$ (+3.9) & $-0.026$ & $-0.005$ & $-0.020$ (-1.7) \\
CLIP-B & $-0.025$ & $-0.044$ & $+0.019$ (+1.7) & $-0.014$ & $-0.010$ & $-0.004$ (-0.4) \\
CLIP-L & $-0.013$ & $-0.049$ & $+0.036$ (+7.2) & $-0.012$ & $-0.012$ & $-0.001$ (-0.1) \\
SigLIP-B & $-0.023$ & $-0.050$ & $+0.027$ (+2.9) & $-0.009$ & $-0.010$ & $+0.001$ (+0.1) \\
\bottomrule
\end{tabular}
\caption{Star-calibrated depth test (excess B, Haar construction). For each backbone the class centroids are compared with a \emph{matched star}: $K$ Gaussian hubs with the real hubs' RMS radius and, around each hub, an isotropic Gaussian cloud with that superclass's real within-cluster RMS spread (same $n$, $d$, $K$ and cluster sizes; 3 star seeds). ``real'' and ``star'' are excesses of $\hat\delta$ over the hub-randomizing null (each cluster kept intact, hubs replaced by a Haar-rotated sample with the hubs' exact spectrum, 10 replicates), which has a small star bias of its own; ``depth'' is their difference (real minus star; negative = more tree-like than a matched star), with $z$ against the combined star, null and estimator spread. On ImageNet only DINOv2-L exceeds its matched star beyond two spreads ($z{=}-2.5$; DINOv2-G $-1.7$, all others $|z|\le1.6$); on CIFAR-100 every backbone is \emph{less} tree-like than its matched star. The excess of \S\ref{sec:form} is therefore accounted for by clustered, star-like organization; residual depth is at most marginal.}
\label{tab:b29-depth}
\end{table}
